% !TeX program = xelatex
% !TeX encoding = UTF-8
% =============================================================================
% main_DIS.tex — Variante DIS (dyslexic-friendly).
%
% Magic comment "% !TeX program = xelatex" sopra: editor moderni (Overleaf,
% TeXstudio, VSCode LaTeX Workshop, TeXworks) compilano automaticamente
% con XeLaTeX, che e' richiesto per il font OpenDyslexic via fontspec.
%
% Il documento usa \verificaFontNormal (sans-serif standard). Solo
% l'\input{esercizi_DIS} viene avvolto in {dyslexicScope}, che applica
% OpenDyslexic + interlinea 1.6 (con auto-fallback a helvet se fontspec
% non e' disponibile o il font non e' installato). Intestazione, ulteriori
% misure, misure dispensative e griglia restano nel font normale.
% =============================================================================
\documentclass[12pt]{article}
\usepackage{{{TEXCOMMON_DIR}}/verifica}
\verificaFontNormal
\verificaSetupFontDis
\verificaCounterFooter{DIS}{{{COPIE_DIS}}}

% G27.tikz.hoist — preamble macros estratti dai TikZ template
\input{tikz_preamble}

\title{ {{TITOLO_VERIFICA}} \\\large [DIS] }
\date{}

\begin{document}
\thispagestyle{firstpage}
\input{{{TEXCOMMON_DIR}}/intestazione}

\begin{dyslexicScope}
\input{{{ESERCIZI_FILE}}}
\end{dyslexicScope}

\input{{{TEXCOMMON_DIR}}/ulteriori_misure}
\input{{{TEXCOMMON_DIR}}/BES_DSA/misure_dispensative}

\input{{{GRIGLIE_DIR}}/{{INDIRIZZO_CODE}}_{{MATERIA_CODE}}{{GRIGLIA_SUFFIX}}}

% Compensazione orale SOTTO la griglia (porting legacy master).
% Inclusa solo se Compensa attivo: TexBuilder risolve {{COMPENSA_OPEN}} a vuoto
% se compensa=true, altrimenti a `%` per commentare l'\input.
{{COMPENSA_OPEN}}\input{{{TEXCOMMON_DIR}}/BES_DSA/compensazione_orale}{{COMPENSA_CLOSE}}
\end{document}
